\scanpage{20}
% L1–4
fixé) d'un $T_i,E_i,B_i\subset E_i$ équivaut en fait
à la donnée de $B_i$ (comme torseur sous
$\pm1$ !) en prenant
\[
T_i=\mathrm{Ker}\,(k^{B_i}\overset{\varepsilon}{-}k),\quad E_i=\varepsilon^{-1}(1)
\]

% L5–10
Donc la donnée \deleted{des} \deleted{\unsure{de la}} \inserted{d'une} famille $(T_i,E_i,B_i)_{i\in I}$
\inserted{($\operatorname{Card} I=n$)} équivaut à
la donnée d'un \unsure{ens orienté|squelette|ens orient.} de $n$-cube.
Montrons que cette donnée permet de
reconstruire le $n$-cube, comme le
produit de la famille (indexée par $I$)
des $1$-cubes ---

% L11–14
On a aussi
\[
\begin{aligned}
T&=\mathrm{Ker}\,(k^B\overset{\overset{\nearrow\,\text{trace}}{\mathrm{tr}}}{-}k^I)\\
E&=\mathrm{tr}^{-1}(1)
\end{aligned}
\]
Si $J\subset I$, on considère
% archive sheet number: 19
